\title{X-LoRA: Mixture of Low-Rank Adapter Experts, a Flexible Framework for Large Language Models with Applications in Protein Mechanics and Molecular Design}

\begin{abstract}
We present a mixture of expert approach for fine-tuning large language models, utilizing a deep layer-wise token-level method grounded in low-rank adaptation (LoRA). Beginning with a collection of pre-trained LoRA adapters, our gating mechanism leverages hidden states to dynamically combine adapted layers, enabling the resulting X-LoRA model to harness diverse capabilities and generate novel deep layer-wise combinations to address tasks. This design draws inspiration from biological concepts of universality and diversity, where neural network components are reused across various hierarchical configurations. Consequently, the X-LoRA model can be seamlessly applied to any existing large language model (LLM) without requiring changes to its underlying architecture. We develop a specialized X-LoRA model offering scientific functionalities such as forward and inverse analysis tasks and improved reasoning abilities, with a focus on biomaterial analysis, protein mechanics, and design. The significance of this work lies in providing accessible, extensible, and adaptable models possessing strong domain expertise and the ability to integrate knowledge across disciplines. Incorporating experts in biology, mathematics, reasoning, bio-inspired materials, mechanics and materials science, chemistry, protein biophysics, mechanics, and quantum mechanics-based molecular properties, we conduct a series of physics-centered case studies. We evaluate knowledge recall, protein mechanics forward and inverse problems, protein design, adversarial agentic modeling including ontological knowledge graph construction, as well as molecular design. The model not only delivers quantitative predictions of protein nanomechanical properties and quantum mechanical molecular characteristics but also interprets the results and accurately identifies plausible mechanisms that elucidate distinct molecular behaviors.
\end{abstract}

\section{Introduction}
Large language models (LLMs)~\cite{Vaswani2017AttentionNeedc,Touvron2023LlamaModels,OpenAI2023GPT-4Report,Chowdhery2022PaLM:Pathways,Jiang2023Mistral7B,Gunasekar2023TextbooksNeed} have become highly popular, including the creation of specialized models that excel in specific tasks, reasoning, or scientific fields~\cite{Bubeck2023SparksGPT-4,Buehler2023MechGPTModalities,Nejjar2023LLMsAnalysis,Buehler2023GenerativeDesign,Luu2023BioinspiredLLM:Materials,Luu2023GenerativeSolvents,Buehler2023MeLMProblemsb,Ge2023OpenAGI:Experts}. Nevertheless, training such models can be expensive, particularly when a broad range of capabilities is required. Techniques like low-rank adapters (LoRA) \cite{hu2021lora} have been introduced as a more resource-efficient alternative, although these adaptations generally concentrate on more specialized knowledge areas. The core idea behind LoRA is to use low-rank matrices added to the original full-sized matrices, designating these low-rank components as the sole trainable parts of the model. Because only these adapter layers are trainable, the models tend to retain the knowledge obtained during the base model’s pre-training while becoming better suited for specific tasks, with the added advantage of computational efficiency. Given that training LoRA adapters is relatively inexpensive, this approach allows for the development of many specialized models.

Although LoRA adapters offer an efficient path to creating specialized model capabilities, the challenge of combining multiple such adapters into enhanced models that integrate and potentially improve capabilities remains unresolved. In fact, combining LLMs has attracted growing interest as a research topic~\cite{Kim2023SOLARUp-Scaling}. Experimental methods, such as spherically interpolated model parameters (SLERP) and other approaches~\cite{Arcee-ai/mergekit:Models.}, have demonstrated promising results and suggest the potential for mixing models to achieve novel capabilities. Mixed models produced via these methods are typically computed {\it a priori} following specific algorithms and, although they have shown encouraging performance on certain benchmarks, further research is needed to systematically develop mixed models and understand how particular features arise.

We propose that dynamically mixing parameters represents a more general methodology, as it enables the combined model to investigate novel parameter combinations during inference.

Similar methodologies have employed mixture of experts (MoE) frameworks~\cite{Jacobs1991AdaptiveExperts,Hampshire1992TheRecognition,Jordan1994HierarchicalAlgorithm}, including the recently introduced Mixtral LLM~\cite{Jiang2024MixtralExperts}. Nonetheless, these approaches often incur high computational costs, imposing substantial demands even at inference time. In contrast, we introduce a computationally efficient mixture-of-experts method that utilizes a collection of distinct LoRA adapters, each of which can be trained independently with ease. Our technique embodies a highly granular and adaptable strategy, inspired by a biological model that reutilizes neural network modules to build a hierarchy of interacting models, spanning from foundational models to agentic systems (Fig.~\ref{fig:hierarchical_concept}). To quantitatively demonstrate its effectiveness, we focus our investigation on leveraging this strategy within science-oriented LLMs, particularly applied here to protein mechanics challenges within the broader context of bio-inspired materials analysis and design.

We present an algorithm designed to dynamically combine adapters at the token-level, offering fine-grained control that permits mixing adapters across individual layers. This method, termed X-LoRA, provides access to a wide variety of capabilities essential for scientific applications involving multimodal language modeling~\cite{Hu2023DeepScience,Buehler2023AMetamaterials,Buehler2022ProteinNetworks}.

\subsection{Fundamental concepts of X-LoRA}

To begin, we briefly revisit the principles underlying the development of LoRA adapters. The core idea of low-rank adaptation (LoRA)~\cite{hu2021lora} posits that weight updates possess a low “intrinsic dimension,” and exploits this by freezing the original weight matrix $W_0 \in \mathbb{R}^{d \times k}$ while restricting modifications to a low-rank decomposition

\begin{equation}
W_0 + \Delta W_0 = W_0 + BA 
\end{equation}

where $B \in \mathbb{R}^{d \times r}$ and $A \in \mathbb{R}^{r \times d}$ with rank $r \ll {\rm min}(d,k)$. During LoRA training~\cite{hu2021lora}, the pretrained weights $W_0$ remain fixed, and only the matrices $A$ and $B$ are updated. The forward pass of a LoRA layer can be written as:

\begin{equation}
    h = W_0x + \Delta W_0x = W_0x + BAx
\end{equation}

In practice, LoRA adapters include a fixed scalar weight $\alpha$, which is applied to the decomposition matrices, resulting in

\begin{equation}
    h = W_0x + BAx \cdot \alpha
\end{equation}

Building upon this, X-LoRA introduces the concept of employing a collection of LoRA adapters, each possessing distinct functionalities, as demonstrated in prior science-focused LLM studies (e.g., \cite{Buehler2023GenerativeDesign}). Rather than statically mixing or integrating models, we propose a dynamic gating mechanism that scales each individual LoRA adapter with granularity over tokens and layers to enable mixing deep within the model. Concretely, for a given LoRA adapter $i$, the original scalar $\alpha_i$ is adjusted by multiplying it by a scaling factor $\lambda_i$, generating the new scaling factor $\alpha^*_i$:

\begin{equation}
    \alpha^*_i = \alpha_i \cdot \lambda_i
\end{equation}

The scaling factor $\lambda_i$ is generated by an X-LoRA scaling head that leverages the model's hidden states; this scaling head constitutes the trainable portion of the X-LoRA architecture. Given that the scaling head can exploit complex encodings, it can be efficiently implemented as a simple feed-forward neural network with one or a few layers (in X-LoRA’s implementation, users can define the architecture of this network).

This scaling operation is performed for every token in the input sequence, thereby providing fine-grained control as each LoRA adapter operating at specific layers within the LLM is individually scaled.

From this point forward, we refer to this method of gating different adapter components as scaling. Additional details about this technique are described in the Materials and Methods section.

\subsection{Paper outline}

The structure of this paper is organized as follows. Initially, we describe the methodology and training framework, including the pertinent datasets, which lead to the creation of an X-LoRA model endowed with specialized abilities in the physical sciences, with a particular emphasis on biomaterials such as proteins. Subsequently, we showcase a collection of experiments where the X-LoRA model is applied to various tasks, including question answering, conversational and agent-based modeling, protein design and analysis, among others. We perform an in-depth examination of the scaling behaviors and validate the approach by comparing it with molecular modeling and other physical data and methodologies.

\section{Results and discussion}

The development process for the X-LoRA model involves the following stages:

\begin{enumerate}
    \item Training the foundational base LLM (accomplished in prior work) \item Training individual adapters to cultivate expertise across several domains (each representing distinct or overlapping areas of knowledge and skills) \item Training the combined X-LoRA model
\end{enumerate}

Our experiments begin with training multiple LoRA adapters. We create nine adapters, each fine-tuned with specific expertise, based on the Zephyr-7B-$\beta$ model~\cite{HuggingFaceH4/zephyr-7b-betaFace}, which itself is built upon the Mistral-7B model~\cite{Jiang2023Mistral7B}:

\begin{enumerate}
    \item Bioinspired materials \item Chain-of-thought (CoT) and reasoning \item Chemistry \item Mathematics \item Physics \item Biology \item Mechanics and materials \item Platypus: Logic and reasoning \item Protein mechanics tasks (generative sequence-to-property and inverse capabilities)
\end{enumerate}

These specific domains of focus and expertise were selected because our ultimate goal is to develop an LLM centered on STEM with a strong emphasis on materials science. Nevertheless, we hypothesize that our methodology can be broadly applied to various fields, particularly for generating modalities that facilitate connections between different knowledge domains~\cite{Buehler2023MechGPTModalities,Buehler2023MeLMProblemsb}.

Following the creation of the LoRA adapters, we proceed to train the X-LoRA model. This training utilizes a subset derived from the combined pool of training samples initially employed for training the individual LoRA adapters (further details can be found in the Materials and Methods section, which includes comprehensive descriptions of the dataset, implementation, and related aspects).

The X-LoRA system serves as an experimental framework to evaluate its effectiveness across multiple scenarios. We perform a range of experiments, such as tackling complex question answering, employing the generative model for protein design, conducting analysis and reasoning tasks, and devising more advanced interaction strategies via adversarial agentic modeling. Throughout these evaluations, we compare the performance of the base model against the X-LoRA model, monitoring how the scaling head leverages various layers from the distinct experts during answer generation or development.

\subsection{Question answering and observed X-LoRA layer-wise scaling weights}

Our initial investigation focuses on the differences between the X-LoRA model and the foundational base model. We pose the following question to both models:

\begin{quote}
\small\texttt{What material behavior (stiffening vs softening) causes supersonic crack speeds in mode II cracks? Think step by step.}
\end{quote}

The base model replies as follows:

\begin{quote}
\small\texttt{Supersonic crack speeds, also known as Rayleigh waves, occur in mode II cracks due to the stiffening material behavior of the material near the crack tip. In mode II fracture, the loading is shear-dominated and the crack propagates perpendicular to the direction of applied stress. Here's a step-by-step explanation:

As the crack tip moves closer to the material, the stresses nearby escalate sharply. This creates a zone of intense stress concentration known as the stress singularity.

Within this stress singularity, the material undergoes elastic deformation that results in an increase in stiffness, i.e., a stiffening effect. This stiffening occurs because the material adjacent to the crack tip becomes increasingly constrained as the crack nears [...]}
\end{quote}

Based on our own expertise~\cite{Buehler2008AtomisticFailure}, this explanation is not only inaccurate but also unnecessarily verbose. By contrast, the X-LoRA model provides the response:

\begin{quote}
\small\texttt{The material behavior that causes supersonic crack speeds in mode II cracks is stiffening.}
\end{quote}

This response is accurate, concise, and clear~\cite{Buehler2008AtomisticFailure}.

The distinction between the base model and the X-LoRA model becomes even more apparent in specialized protein analysis and design tasks. The task presented is:

\begin{quote}
\small\texttt{CalculateForceHistory< M S K S N T Y R M L V L L E D D T K I N K E D E K F L K G K P G K M H E F V D E L I L P F N V D E L D E L N T W F D K F D A E I C I P Y V E I L K E S G M K > }
\end{quote}

The base foundation model replies with:

\begin{quote}
\small\texttt{I'm sorry, but the input you provided is not a valid mathematical expression. It seems to be a sequence of characters that do not constitute a computable function. Please supply a valid mathematical expression for calculation.

Nonetheless, I can mention that the sequence you gave consists of function and variable names that might be part of a larger program or script in a particular programming language. Without additional context, determining the intended computation is not possible.

If you provide further details, such as the programming language and the context of this expression, I might be able to offer more clarity about its meaning and potential calculations.}
\end{quote}

In contrast, the X-LoRA model responds as follows, successfully accomplishing the design task of producing the force-deformation curve related to this protein:

\begin{quote}
\small\texttt{[0.013,0.066,0.078,0.045,0.072,0.126,0.127,0.107,0.119, 
0.074,0.106,0.098,0.103,0.128,
0.077, 0.069,0.131,0.108,0.144,0.174, 0.166,0.183,0.185,0.217,0.175,0.230,0.187,0.224, 0.267,0.301,0.315]}
\end{quote}

Figure~\ref{fig:QA_weights_sample_0} presents outcomes from this question-answering task along with the observed X-LoRA scaling weights. The two prompts, one centered on mechanics and materials (Fig.~\ref{fig:QA_weights_sample_0}a) and the other on protein mechanics (Fig.~\ref{fig:QA_weights_sample_0}b), exhibit distinct activation patterns. Notably, a complex arrangement of scaling values is utilized in each instance, implying that the X-LoRA model leverages a heterogeneous combination of different adapters across layers. Additionally, the heatmap in Fig.~\ref{fig:QA_weights_sample_0}a shows a wider spread of activation among numerous experts and layers, featuring some regions with strong activation. Conversely, the heatmap in Fig.~\ref{fig:QA_weights_sample_0}b displays a more focused activation pattern, with fewer regions exhibiting high values. This observation suggests that the gating function on the right may be more selective in triggering experts.

In Fig.~\ref{fig:QA_weights_sample_0}, the bright line-like regions indicate pathways of strong activation, implying that the scaling function preferentially engages certain experts for particular layers. The observed patterns suggest that the importance is not evenly distributed across all experts and layers; rather, it is relatively sparse and selective. This sparsity aligns with findings in other mixtures-of-experts models, where only a limited subset of experts is activated for any given input to specialize in distinct parts of the problem domain.

By comparing the activation maps across different tasks, both the sparsity and the specific patterns may provide valuable insights into which experts the model depends on for various input types or during different phases of processing within the neural network (the dynamic evolution of these maps during complex task solving will be discussed later).

To further investigate this behavior, we examine a broader range of questions to understand how the X-LoRA model leverages different experts when presented with diverse prompts. Our focus is on determining whether, and how, the X-LoRA model combines multiple adapters for questions that span overlapping skills or domains. Fig.~\ref{fig:QA_weights} displays the outcomes from question answering along with the observed X-LoRA scaling weights, corresponding to the queries listed in Table~\ref{tab:table_qa}. Table~\ref{tab:table_qa} includes both the responses from X-LoRA and from the base model, Zephyr-7B-$\beta$, highlighting significant differences between the two.

To clearly illustrate these differences, we highlight incorrect or questionable answers in \redhighlight{red highlight}. The comparison demonstrates that X-LoRA delivers substantially higher accuracy along with more concise responses.

As anticipated, we detect complex adapter mixing with frequent activation of multiple dominant LoRA experts. For example, Fig.~\ref{fig:QA_weights}d shows that a question about the mechanics of pine cones triggers strong activation of the mechanics/materials adapter combined with the CoT/reasoning and bioinspired adapters. Although X-LoRA selects the correct letter, the answer is not fully accurate since pine cones open under dry, not wet, conditions. However, a follow-up question probing whether release occurs in dry or wet conditions yields the correct answer.

In a protein-related question, Fig.~\ref{fig:QA_weights}f reveals adapter mixing in response to an inquiry about the mechanics of Bombyx mori silk, activating a blend of the bioinspired, mechanics/materials, and biology adapters. For highly specialized tasks like protein design (Fig.~\ref{fig:QA_weights}a), we observe predominantly singular activation of a specific expert.

A more detailed examination of the heatmap of scalings shown in Fig.~\ref{fig:QA_weights}a reveals that besides the clearly significant role of the protein mechanics expert, there are multiple bands where certain experts receive considerable importance across several layers. This pattern is especially prominent for the mechanics/materials and reasoning experts. These bands are intermittent and exhibit a fluctuating pattern, which could imply that these experts are particularly effective under specific conditions or for certain features. Notably, the reasoning expert (positioned just to the left of the protein mechanics expert) is distinguished by several layers where it attains the highest activation, particularly in the upper layers (approximately layers 100 to 140). We infer that this might indicate an important contribution to the model’s final decision-making or output, though further investigation is required to confirm this. Additionally, there is a more focused yet strong activation for the physics and biology experts at particular layers, possibly reflecting a specialization in handling features or representations unique to those layers.

Moreover, Fig.~\ref{fig:QA_token_history} presents a token-level history aggregated over all layers’ scalings. This analysis highlights token histories for two cases: Fig.~\ref{fig:QA_token_history}a relates to a task involving pine cone seed release, while Fig.~\ref{fig:QA_token_history}b depicts outcomes from a sequence of protein mechanics tasks. The protein mechanics task entailed multiple prompt-answer interactions: initially two property calculations followed by a generative task aimed at designing a new protein near the end. Although the engagement of experts remains generally consistent, noticeable variations occur in the usage of the LoRA experts throughout the process. Specifically, the protein mechanics expert is heavily engaged in the final stages of generation when the new protein design is taking place.

Table~\ref{tab:table_qa} illustrates various use cases spanning different domains and compares performance relative to the base model operating independently. We explore several of these examples in greater detail.

For example, we present a question that combines protein science and biology with logical reasoning, inquiring about protein expression patterns. X-LoRA provides the correct answer. The base model's response, however, contains some errors and logical flaws in its analysis of the marker protein expressions in each cell type. The base model accurately recognizes that fibroblasts express Protein B and do not express Protein A due to the mutual exclusivity of Protein A and Protein B. Nevertheless, it fails to explicitly state that fibroblasts express Protein C. Since there is no rule preventing the co-expression of Protein C with Protein B, and fibroblasts have no restriction on Protein C, it logically follows that fibroblasts express Protein C as well. Regarding Neurons, the base model begins correctly by noting that neurons express Protein A and not Protein B based on the mutual exclusivity rule. Yet, the discussion about neurons potentially expressing Protein C or both Protein C and Protein B is confusing and contains an error. Because neurons express Protein A and not Protein B (as indicated), the only remaining possibility, given the rules, is that they express Protein C. There is no circumstance in which neurons could express Protein B, contrary to the initial analysis. Consequently, neurons express Protein A and Protein C, not Protein B. For Epithelial cells, the base model correctly points out that epithelial cells do not express Protein A but do express both Protein B and Protein C. However, the explanation is convoluted and includes incorrect assertions regarding the mutual exclusivity of Protein A and Protein C, which was not stated in the original information. The straightforward reasoning is that epithelial cells do not express Protein A but must express the two other proteins, Protein B and Protein C, since the original instructions only specify that Protein A and Protein B cannot be co-expressed.

X-LoRA, on the other hand, accurately identifies the marker proteins expressed by each cell type based on the given observations and logical reasoning. For fibroblasts, X-LoRA correctly infers that fibroblasts do not express Protein A because the presence of Protein A excludes the expression of Protein B. Since fibroblasts do express Protein B, and considering that each cell type expresses a distinct combination of three marker proteins (an incorrect assumption introduced, as the original statement did not specify the number of proteins each cell type expresses but only that they express a unique combination), the conclusion that fibroblasts also express Protein C is valid given the information provided (excluding the incorrect detail about the number of proteins, which appears to be a misunderstanding). Regarding neurons, the model accurately determines that they express Protein A and Protein C. The logic is sound, noting that Protein B cannot co-express with Protein A, and since neurons express Protein A and one other marker protein (not Protein B), they must express Protein C. Lastly, for epithelial cells, X-LoRA correctly concludes that these cells express Protein B and Protein C. This is because epithelial cells lack Protein A expression, and by elimination, since they express two marker proteins, these must be Protein B and Protein C.

In the query concerning pine cone release under humid versus dry conditions, the base model provides a convoluted and factually inaccurate response. X-LoRA offers a precise answer accompanied by clear reasoning. Similar examples include the question about the exceptional mechanical properties of Bombyx mori silk, where the base model mistakenly claims that silk fibers conduct electricity. The base model also fails on several other domain-specific questions, such as the one regarding the role of hyperelasticity in maximum fracture speeds, where X-LoRA delivers a concise and correct answer while the base model is unable to respond accurately. In a question about the strong adhesion of gecko feet, X-LoRA correctly emphasizes the significance of setae and nanoscale effects, whereas the base model incorrectly attributes this to specialized glue proteins. These and other examples presented in the table and elsewhere in the paper provide compelling evidence of X-LoRA’s superior performance compared to the base model.

Figure~\ref{fig:perf_comp_bioinspiredexam} presents a comparison of knowledge recall evaluation experiments utilizing the bio-inspired knowledge test set introduced in \cite{Luu2023BioinspiredLLM:Materials}. It is evident that X-LoRA delivers the best performance, despite being a significantly smaller model than both the BioinspiredLLM and the Llama-BioLLM models (7B parameters for X-LoRA versus 13B parameters in the others). The figure also includes several comparisons with other recently published LLMs, specifically Orca-13B and Llama-13b-chat.

Additionally, the plot shows a comparison between X-LoRA and its base model, Zephyr-7B-$\beta$, demonstrating that X-LoRA achieves notably improved performance. These results suggest that strong performance is attainable with X-LoRA even when smaller base models are employed.

Figure~\ref{fig:image_synth} illustrates an application in image synthesis, where X-LoRA is used to generate prompts for other models, such as DALL-E 3, as demonstrated here.

\begin{quote}
\small\texttt{Create a prompt that I can use for image generation that accurately describes the microstructure of a hierarchical composite.

Explicitly describe the geometry.}
\end{quote}

The comparison shown in Figure~\ref{fig:image_synth}(b)-(c) highlights that the prompt generated by X-LoRA is more concise, which leads to a substantially more accurate image produced by DALL-E 3~\cite{DALLECard}. Importantly, the prompt created by Zephyr-7B-$\beta$ does not faithfully represent the given instruction and adds various other design elements, such as fibers and nanoparticles, which were not mentioned in the original prompt focused solely on the microstructure of a hierarchical composite. The X-LoRA prompt is more succinct and targeted, resulting in a significantly better output.

\subsection{Generative solutions to protein analysis}

We now turn to specific generative and analytical tasks related to proteins. These abilities arise in the X-LoRA model due to the protein mechanics adapter and are further enhanced by integrating the model’s knowledge of biology, bio-inspired materials, reasoning, and logical deduction. This section focuses on the analysis of protein-related tasks independently, evaluating the quantitative performance of this functionality.

Fig.~\ref{fig:protein_force_history_prediction} presents outcomes from the protein task focused on predicting force-deformation behaviors based on sequences. Fig.~\ref{fig:protein_force_energy_prediction} displays the overall accuracy in forecasting unfolding force and unfolding energy from the sequence, comparing the true values with the model's predictions (yielding $R_2=0.85$). In general, the model demonstrates excellent performance and, as will be examined in the next section, also performs very well for the inverse problem of design and cycle consistent evaluation\cite{Lu2023GenerativePropertiesb}. However, unlike previous studies that relied on highly specialized GPT-type models for these tasks, our X-LoRA model integrates these specialized abilities along with a variety of other deep specializations obtained through multiple fine-tuned adaptations.

\subsection{Protein design and analysis}

We now extend the range of test cases to include more comprehensive applications that span several capability domains. Specifically, the experiments presented here illustrate how the X-LoRA model leverages a diverse set of skills and how agentic modeling can facilitate even more complex and thorough responses. We start by discussing a protein design task, where the model is tasked with creating a novel protein to achieve a specific target force-deformation response. These force-deformation behaviors have been investigated in prior work~\cite{Ni2023ForceGen:Model} and represent measurements of force versus deformation as the protein is stretched from both ends (mirroring a classical protein unfolding experiment~\cite{Lu1998UnfoldingSimulation.}).

Fig.~\ref{fig:design_protein} illustrates the application of the generative protein task. We design a protein with a targeted force-deformation response (using training data and boundary conditions from the original molecular dynamics (MD) simulations of protein pulling, see~\cite{Ni2023ForceGen:Model}), and then evaluate the predicted sequence to determine if it achieves the specified target properties (Fig.~\ref{fig:design_protein}a). This cycle-consistent approach is crucial for verifying whether the inverse design capabilities align with the forward predictions.

In Fig.~\ref{fig:design_protein}b, the predicted sequence \texttt{MSKSNTYRMLVLLEDDTKINKEDEKFLKGKPGKMHEFVDELILPFNVDELDELNTWFDKFDAEICIPYVEILKESGMK} is shown. Additionally, Fig.~\ref{fig:design_protein}c presents a folded protein model generated by AlphaFold 2~\cite{Jumper2021HighlyAlphaFoldb}. To evaluate the relationship of this designed protein to known sequences, Fig.~\ref{fig:design_protein}d provides a Basic Local Alignment Search Tool (BLAST) Tree~\cite{Altschul1990BasicTool} comparison. Notably, the designed protein shares certain connections with hypothetical proteins previously studied but not yet structurally characterized or identified.

By employing agentic modeling and integrating multiple X-LoRA models as components in a multi-agent framework, we further investigate the predicted sequence through our X-LoRA model. This leverages our model’s enhanced capacity to not only interpret protein mechanics but also to combine this insight with other knowledge domains. Using the predicted sequence as input, we request a more detailed analysis and suggestions for sequence variations. We begin with the following prompt:

\begin{quote}
\small\texttt{I have identified this amino acid sequence: MSKSNTYRMLVLLEDDTKINKEDEKFLKGKPGKMHEFVDEL ILPFNVDELDELNTWFDKFDAEICIPYVEILKESGMK. 

How can I use it? 

Provide logical steps, discuss specifically the engineering of the problem. }
\end{quote}

Multiple agentic interactions are utilized to generate the solution, particularly employing adversarial strategies (for details, see Materials and Methods) to encourage the model to explore more profound and complex aspects of the discussed concepts. The following summary of this dialogue is produced by the X-LoRA model (the complete conversation is available in the Supplementary Material, SI Conv Example 1):

\begin{quote}
\small\texttt{The discussion centers on the potential uses of a particular amino acid sequence in the field of protein engineering. This sequence features both hydrophobic and hydrophilic residues, implying it may form a stable conformation and engage with other molecules. To design a protein incorporating this sequence, computational approaches can predict its secondary and tertiary structures; however, challenges include the complexity of protein folding and the high computational demands. Enhancements can be achieved by creating more efficient algorithms and integrating new experimental findings. It is essential to validate computational models to guarantee their accuracy and practical relevance. Methods such as cross-validation, regularization, early stopping, ensemble techniques, and transfer learning are effective for avoiding overfitting and underfitting. Additionally, interpreting and explaining models aids in guiding practical implementations.}
\end{quote}

We also request the model to generate a summary of main points, yielding:

\begin{quote}
\small\texttt{\begin{itemize}
    \item The amino acid sequence MSKSNTYRMLVLLEDDTKINKEDEKFLKGKPGKMHEFVDELI LPFNVDELDELNTWFDKFDAEICIPYVEILKESGMK holds promise for protein engineering because of its combination of hydrophobic and hydrophilic residues. Computational techniques can predict the protein’s secondary and tertiary structures, though they face limitations like protein folding complexity and computational expenses.
    \item Validating computational models is vital to confirm their precision and applicability. Strategies such as cross-validation, regularization, early stopping, ensemble methods, and transfer learning help mitigate overfitting and underfitting.
    \item Model interpretation and explanation are also valuable for guiding practical uses.
\end{itemize}
}
\end{quote}

Ultimately, the key insight is:

\begin{quote}
\small\texttt{The key insight is that the amino acid sequence MSKSNTYRMLVLLEDDTKINKEDEKFLKGKPGKMHEFVDELILPFNVDELDELNTWFDKFDAEICIPYVEILKESGMK holds promise for protein engineering because it contains a combination of hydrophobic and hydrophilic residues. To design a protein incorporating this sequence, computational approaches can be utilized to predict its secondary and tertiary structures; however, challenges include the complexity of protein folding and the high computational demands. It is essential to validate computational models to guarantee their accuracy and practical applicability. Methods such as cross-validation, regularization, early stopping, ensemble techniques, and transfer learning can help mitigate overfitting and underfitting. Moreover, interpreting and explaining the models contributes valuable insights for their practical use.}
\end{quote}

Further analysis can delve into detailed sequence engineering strategies. For this purpose, the model is prompted to suggest specific design modifications aimed at enhancing the protein’s stability:

\begin{quote}
\small\texttt{What changes could be made to the following protein sequence to improve its stability: MSKSNTYRMLVLLEDDTKINKEDEKFLKGKPGKMHEFVDELILPFNVDELDELNTWFDKFDAEICIPYVEILKESGMK?}
\end{quote}

A summary of the dialogue is as follows (see Supplementary Material, SI Conv Example 2, for the full exchange):

\begin{quote}
\small\texttt{The discussion centers on strategies to increase protein stability by incorporating hydrophobic amino acids and disulfide bonds. The biologist notes that hydrophobic residues typically enhance protein stability, and that disulfide bonds can also contribute to structural reinforcement. The engineer requests concrete examples of hydrophobic amino acids and the recommended number of disulfide bonds to include. The biologist cites examples of proteins—such as insulin, bovine pancreatic trypsin inhibitor, and green fluorescent protein—that have been successfully stabilized using these approaches. The conversation also addresses potential disadvantages of adding hydrophobic residues or disulfide bonds, including altered protein function, added complexity, and the risk of aggregation. The biologist recommends overcoming these issues through careful amino acid selection, computational modeling, and experimental validation.}
\end{quote}

The summary of main points:

\begin{quote}
\small\texttt{\begin{itemize}
\item Stability of a protein sequence can be enhanced by incorporating hydrophobic amino acids and disulfide bonds.

Hydrophobic residues generally contribute to protein stability, and disulfide bonds also serve to reinforce stability.

Proteins such as insulin, bovine pancreatic trypsin inhibitor, and green fluorescent protein have been successfully stabilized through these modifications.
\item Possible challenges or disadvantages of adding hydrophobic amino acids or disulfide bonds include altered protein function, increased structural complexity, and risk of aggregation.
\item These challenges can be mitigated by careful amino acid selection, computational modeling, and empirical testing.
\end{itemize}
}
\end{quote}

In conclusion, the key insight is:

\begin{quote}
\small\texttt{The most crucial point from the discussion is that to improve the stability of the protein sequence MSKSNTYRMLVLLEDDTKINKEDEKFLKGKPGKMHEFVDELILPFNVDELDELNTWFDKFDAEICIPYVEILKESGMK, introducing hydrophobic amino acids and disulfide bonds is an effective strategy. This directly addresses the initial question.}
\end{quote}

From the viewpoint of protein engineering, these recommendations are logical because they can lead to designs with enhanced stability. Detailed examples and methodologies have been formulated that researchers can further investigate.

As another illustration, consider the following three sequences. We initially ask the model to evaluate the strength of each sequence and then to analyze the outcomes: identify the strongest candidate and provide reasoning for why this protein is likely the most stable. Additionally, the model is requested to describe the steps to synthesize it in the laboratory. The three sequences (manually derived by altering one sequence from the test set) are:

\begin{quote}
\small\texttt{
\begin{itemize}
    \item LNTWFDKFDAEICIPYVEILKESGMK \item AITWFDKFDAEICIPYVEALKIAAMK \item AAAAAAAAKFDAAEICIIAAAAAAAAKIAAMK
\end{itemize}
}
\end{quote}

This scenario assesses the model’s ability to integrate multiple skills and adapt dynamically as the conversation unfolds. Specifically, it tests the capacity to switch flexibly between areas such as protein structural analysis, logical reasoning, hypothesis generation, and more.

The conversation is illustrated in Figure~\ref{fig:conversation_evolve}, along with a visualization of how the scalings evolve throughout the dialogue (the analysis presents both the deep layer-wise adapter scaling and an aggregated analysis to provide an indicator of the effective utilization of the various experts). The findings clearly demonstrate that the protein task is initially highly relevant, progressively shifting towards a stronger emphasis on the bioinspired adapter and the biology adapter. By the conclusion of the conversation, the biology adapter becomes the most dominant.

To evaluate the predictions and assess the quality of the responses, the three proteins examined in this task are shown in Figure~\ref{fig:protein_visuals} (proteins folded using Alpha Fold 2~\cite{Jumper2021HighlyAlphaFoldb}, via ColabFold \cite{Mirdita2022ColabFold:All}). We observe that the most stable structure, located in the middle, exhibits the most organized secondary structure, which aligns with the idea that it produces the highest unfolding force, as predicted by the X-LoRA model in Figure~\ref{fig:conversation_evolve}. The other two structures are less organized; the first one is primarily helical with some turns, and the last one is partially unstructured and displays a kink in its geometry. These characteristics account for the lowest predicted unfolding force as estimated by the model.

By examining the visual depictions of the folded configurations, we verify that the central structure, which is the most stable, exhibits the most well-organized secondary structure. This observation aligns with the idea that it produces the greatest unfolding force as forecasted by the model. We observe that the other two structures are less ordered: the first one is predominantly helical with some turns, while the last one is partly unstructured and contains a bend in its geometry. We propose that these characteristics probably account for the lowest unfolding force detected, consistent with the predictions made by the X-LoRA model.

In a subsequent example, we concentrate on the energetic characteristics of proteins, specifically analyzing the unfolding energy as the protein is extended from its initial conformation to a fully unfolded state~\cite{Ni2023ForceGen:Model}. We again study three sequences (the first two being the same as those previously examined, but the third comprises a longer chain, combining the second sequence at the beginning and end with the first sequence in the middle):

\begin{quote}
\small\texttt{\begin{itemize}
    \item LNTWFDKFDAEICIPYVEILKESGMK \item AITWFDKFDAEICIPYVEALKIAAMK \item AITWFDKFDAEICIPYVEALKIAAMKLNTWFDKFDAEICIPYVEILKESGMKAITWFDKFDAEICIPYVEALKIAAMK
\end{itemize}
}
\end{quote}

Our initial step is to have the model calculate the unfolding energy for each sequence, followed by reasoning about the outcomes: Considering the amino acid sequences, we request an explanation for why this particular protein is probably the most stable. Finally, we inquire about potential functions that the most stable protein might possess.

Figure~\ref{fig:MJB_20} presents a transcript of a dialogue that engages multiple distinct experts and synthesizes knowledge across these domains, demonstrated here through the analysis of the unfolding energy of three proteins. By the conclusion of the conversation, the CoT/reasoning adapter emerges as the dominant element, as the X-LoRA model addresses a question requiring reasoning over the findings to predict probable protein structural characteristics and possible functions. Specifically, the mechanistic inquiry into the basis of protein stability includes the development of hydrophobic interactions among nonpolar amino acids as well as hydrogen bonds between polar amino acids. The model forecasts that several hydrophobic amino acids (I, F, W, Y) alongside hydrogen bonding amino acids (D, E, K) play a role in maintaining protein stability.

Additionally, the model predicts that cysteine (C) residues present in the sequence would promote disulfide bond formation, which further reinforces the protein’s structure. Notably, the critical predictions concerning the structural traits explaining the high stability are accurately identified, as validated by the structural analyses depicted in Fig.~\ref{fig:MJB_21}. This figure’s analysis corroborates the presence of all these features. It also clarifies why the other two proteins, consisting of much shorter segments with variable helical domain stability, exhibit significantly lower unfolding energies.

\subsection{Adversarial agentic modeling to connect distinct scholarly disciplines and knowledge yielding ontological knowledge graph generation}

Given that our X-LoRA model possesses extensive cross-domain expertise, we leverage it to explore links among diverse concepts, knowledge repositories, and specialized fields. In one experiment, we pose two queries to the model, each constructed similarly but with a different focus. Subsequently, we extract triplets to create an ontological knowledge graph, aiming to organize the responses into a more structured format~\cite{Giesa2012CategoryDesign,Spivak2011ReoccurringAnalogies}. The resulting graph offers a cohesive synthesis of the generated insights and illustrates the relationships among concepts in an interpretable and mechanistic way.

We employ two adversarial agents to address the question. Agent 1 assumes the role of a "critical physicist," while Agent 2 acts as a "philosopher." Agent 1 initiates the inquiry and is responsible for probing and generating further questions in response to each reply from Agent 2. As their dialogue progresses, it delves deeper into the topic and reveals multiple facets. The question is deliberately crafted to span various disciplinary domains in order to evaluate the model's capacity to synthesize ideas and concepts.

The question posed is:

\begin{quote}
\small\texttt{Consider an abstract representation of a classical musical composition as a graph of interacting notes, chords, melodies, rhythms and so on.

What happens if we would apply pressure, up to a point of failure, on this abstract concept of music?

Provide logical steps, mathematical reasoning, and discuss specifically the physics of the problem.}
\end{quote}

The complete conversation is provided in the Supplementary Material (see, SI Conv Example 3). For brevity, we present here only the summary and the principal conclusion, which capture the core aspects:

\begin{quote}
\small\texttt{The dialogue explores the effect of applying pressure to an abstract representation of a classical musical composition, treated as a physical system composed of interacting notes, chords, melodies, and rhythms. Pressure can be quantified in terms of stress, and failure occurs when stress reaches a critical threshold. To validate the mathematical model’s predictions, a physical analog of the abstract musical structure can be utilized, allowing experimental results to be compared with theoretical forecasts. The model’s limitations include oversimplified interaction representations, static nature, uniform distribution assumptions, and linear stress-pressure relationships. Enhancing the model’s precision and applicability requires integrating more detailed and realistic features, dynamic behaviors, non-uniform distributions, and nonlinear effects.

[...]

The key insight is that applying pressure to the abstract representation of classical music can be interpreted as altering the force field among the interacting musical elements (notes, chords, melodies, rhythms). Increasing pressure raises the system’s internal stress, and at a critical point, the system fails and collapses under its own structural load. This explanation addresses the original question by providing a physical rationale for how pressure influences the abstract concept of music and how such influence can be quantified through stress.}
\end{quote}

Although the generated text offers valuable insights, constructing an ontological representation of the data can enhance interpretability in analysis~\cite{Buehler2023MechGPTModalities}. For graph generation, we merge the entirety of the two-agent conversation and build a knowledge graph. An example graph is illustrated in Figure~\ref{fig:graph}. The resulting graph is partitioned into communities employing the Girvan-Newman algorithm~\cite{Girvan2002CommunityNetworks}, yielding a total of 12 distinct communities. The figure depicts both the complete graph (Figure~\ref{fig:graph}a) and zoomed-in views with node labels (Figure~\ref{fig:graph}b-c).

\subsection{Creation of X-LoRA-Gemma Incorporating Combined Protein, Chemical, Bio-Inspired, and Materials Mechanics Functions}
\label{gemma-section}

To demonstrate that the proposed method is effective with different foundational models having varied architectures, we trained an additional X-LoRA model, this time built upon the Gemma-7B-it model~\cite{Google/gemma-7b-itFace}. For this model, we utilized four LoRA adapters, which are described as follows:

\begin{enumerate}
    \item Bioinspired materials \item Mechanics and materials \item Protein mechanics tasks (including generative sequence-to-property and inverse design capabilities) \item Quantum-mechanics based molecular properties QM9 (incorporating generative SMILES-to-property and inverse design functionalities, refer to Table~\ref{tab:table_QM9} for definitions and a summary of all properties computed in this dataset)~\cite{Ruddigkeit2012EnumerationGDB-17,Ramakrishnan2015ElectronicSpace}
\end{enumerate}

The development of the X-LoRA-Gemma model follows a similar approach as the original X-LoRA model but is intended to include a distinct set of functionalities, particularly adding the capacity to predict twelve quantum mechanical properties and to design molecules targeting these twelve quantum mechanical properties. This new adapter is trained in conjunction with the previously mentioned ones, using the same datasets from the earlier experiments. Note that while most of the paper centers on the Zephyr-based X-LoRA model, in this section we distinguish the models by referring to the original X-LoRA as X-LoRA-Zephyr and the Gemma-based version as X-LoRA-Gemma.

Figure~\ref{fig:protein_force_energy_prediction} illustrates the overall performance of X-LoRA, showing predictions of unfolding force and unfolding energy from sequence data ($R_2 = 0.87$), as well as the prediction accuracy for quantum mechanical properties of molecules derived from the QM9 dataset ($R_2 = 0.96$). We observe excellent results for both tasks, with slightly improved performance compared to the original X-LoRA model on the protein design task.

Next, we highlight X-LoRA-Gemma’s capabilities through a straightforward multi-agent design experiment. We select a small organic molecule, \texttt{CCCC12CC1C(=O)N2} (sourced from the QM9 test set). This molecule features a bicyclic structure with a ketone (=O) and an amine (N) incorporated in the ring, as visualized in Figure~\ref{fig:QM9_design}(a). We begin by asking the X-LoRA-Gemma model:

\begin{quote}
\small\texttt{Here is a molecule in SMILES code: CCCC12CC1C(=O)N2

What modifications should be made to render this molecule more suitable for polymer design? }
\end{quote}

The response is:

\begin{quote}
\small\texttt{To enhance the utility of this molecule for polymer development, you might:
\begin{enumerate}
\item  Add a reactive functional group, such as an amine or hydroxyl group, to enable further modifications and crosslinking.
\item  Extend the carbon backbone to be longer and more flexible, increasing the molecule's capacity to form polymers with varied structures and characteristics.
\item  Attach a side chain containing a functional group, like a carboxyl group, to facilitate covalent bonding with other molecules and polymers.
\item  Integrate a cyclic component into the molecule, such as a lactam ring, to produce a polymer that is more rigid and stable.
\item  Incorporate a diverse set of side chains to generate polymers exhibiting a range of properties and functionalities.
\item  Vary the number of carbons in the backbone to enable the formation of polymers with different molecular weights and attributes.
\item  Add a crosslinking site within the molecule to form a polymer that is more stable and less reactive.
\item  Introduce a repeating unit in the structure to create polymers with specific molecular weights and features.
\item  Combine several of these approaches to develop a versatile and functional polymer suitable for multiple applications.
\end{enumerate}}
\end{quote}

Next, we concentrate on the initial recommendation.

\begin{quote}
\small\texttt{How can I enhance the overall reactivity of a molecule? I intend to utilize this molecule to construct a polymer material.

The molecule under consideration, represented in SMILES notation, is: CCCC12CC1C(=O)N2

I am able to adjust the following properties:

Dipole moment, mu Polarizability, alpha Highest Occupied Molecular Orbital energy, homo Lowest Unoccupied Molecular Orbital energy, lumo HOMO-LUMO gap: The energy difference between HOMO and LUMO, gap Electronic spatial extent Zero-point vibrational energy Heat capacity at constant volume Internal energy at 0 K Internal energy at 298.15 K Enthalpy at 298.15 K Free energy at 298.15 K What are the top THREE properties I should adjust, and should I increase or decrease them?}
\end{quote}

The reply from X-LoRA-Gemma is:

\begin{quote}
\small\texttt{\begin{itemize}
    \item  HOMO-LUMO gap (decrease) \item  Dipole moment (increase) \item  Polarizability (increase)
\end{itemize}
Adjusting these properties can enhance the overall reactivity of the molecule and render it more appropriate for constructing polymer materials. }
\end{quote}

We leverage the generative and predictive functionalities of X-LoRA-Gemma to engineer these specific modifications in molecular properties.

For the optimization, we employ a dual-agent framework consisting of two instances of X-LoRA-Gemma. One agent performs the generative task \texttt{GenerateMolecularProperties<...>} to propose candidate molecules. The other agent then conducts a computation using \texttt{CalculateMolecularProperties<...>} to evaluate the discrepancy between the desired and the predicted properties. This iterative procedure continues either until the maximum number of iterations is reached or until the mean-squared error falls below a predefined threshold.

Figure~\ref{fig:QM9_design} presents the outcomes of this entire approach. Panel (a) in Figure~\ref{fig:QM9_design} illustrates the target molecular design and the extent of its deviation from the original molecule's properties. Panel (b) displays the top 20 candidate molecules generated by the algorithm, accompanied by the mean-squared error distribution (panel (c)) and a graph of mean-squared error progression across the generated candidates (panel (d)). The optimal design identified is \texttt{CC(C)C1=CC(=O)NC=C1}, which closely approximates the target, as seen in panel (e). Finally, panel (f) in Figure~\ref{fig:QM9_design} shows the refined 3D molecular model, where we apply UFF~\cite{Rappe1992UFFSimulations} for 3D structural optimization.

To evaluate the molecule's reactivity further, we calculate the Gasteiger charges~\cite{Gasteiger1980IterativeCharges} (see Figure~\ref{fig:QM9_design}(f) for a visual summary). The molecule includes an aromatic ring featuring a ketone group (C=O), which is more electrophilic relative to carbon-carbon or carbon-hydrogen bonds, thus prone to nucleophilic attack. The aromatic ring combined with nitrogen and the ketone group forms a pyridone ring structure. This structure exhibits enhanced reactivity compared to a simple aromatic ring because the heteroatom (nitrogen) and the ketone introduce reactive centers, such as the carbonyl carbon prone to nucleophilic addition and potential sites for electrophilic substitution within the ring. The nitrogen atom in the ring, as part of the heteroaromatic framework, influences aromaticity and reactivity by modifying electron density and its distribution throughout the ring. Consequently, this nitrogen-containing ring serves as a locus for chemical reactions, especially those involving electrophilic attack, owing to the electron-donating effect of the nitrogen atom.

The carbon atom within the ring, adjacent to both an oxygen and a nitrogen atom, carries a partial positive charge (+0.25), indicating that it is somewhat deficient in electrons. This is because oxygen, being more electronegative, attracts electron density toward itself, thereby reducing the electron density on the carbon. This influence is partially offset by the neighboring nitrogen, which is less electronegative than oxygen but more so than carbon. Nonetheless, the presence of this partial positive charge signifies that the carbon atom is a probable site for nucleophilic attack, where an electron-rich nucleophile would be drawn to this electron-poor carbon.

The nitrogen atom bears a partial negative charge (-0.33), signifying that it is relatively electron-rich compared to its typical state. This can be attributed to its lone pair of electrons and the influence of the aromatic ring's electron system. In aromatic heterocycles, the nitrogen lone pair can be involved in the delocalized $\pi$-electron network, though the extent depends on the ring’s structure and the electronegativities of neighboring atoms. This partial negative charge implies that nitrogen is more nucleophilic, with an increased tendency to donate electrons, either by forming bonds with electrophiles or through protonation reactions. Additional quantum mechanical studies are necessary to explore this further.

The carbon atom carrying the partial positive charge represents a reactive site for nucleophilic attacks. Other molecules may interact at this position through mechanisms in which nucleophiles target electron-deficient carbons.

The hydrogen atom bonded to the nitrogen, having a partial positive charge of 0.17, is a strong candidate for forming hydrogen bonds with other molecules or atoms. Hydrogen bonds are a form of dipole-dipole interaction occurring when a hydrogen atom, covalently bonded to a highly electronegative element (such as nitrogen, oxygen, or fluorine), interacts with another electronegative atom that possesses a lone electron pair. This feature enables the molecule to act as a hydrogen bond donor, a crucial aspect influencing the molecule’s solubility in water and other polar solvents, as well as its behavior in biological environments.

The ketone group adjacent to the nitrogen-containing aromatic ring imparts distinct reactivity that can be leveraged in the synthesis of polymers~\cite{McQuarrieSimonPhysicalChemistry,CareySundbergAdvancedOrganicChemistry,Varghese2022BeyondApplications,Varghese2022BeyondApplications}. This compound is classified as a lactam due to the cyclic amide formed by the ketone and nitrogen within the ring. Lactams play a significant role in polymer chemistry, especially in producing polyamides, a polymer category widely utilized in various material applications. The ketone functionality is capable of undergoing multiple chemical reactions valuable in polymer chemistry, including nucleophilic addition. Such reactions can be employed to elongate polymer chains or introduce side groups that modify polymer characteristics. The nitrogen atom in the aromatic ring notably influences the polymer's electronic features and can engage in hydrogen bonding, affecting melting temperature, solubility, and mechanical behavior. N-heteroaromatic compounds also exhibit the ability to coordinate with metal ions, enabling potential uses in catalysis or in materials with electronic or photonic functionalities. Furthermore, this molecule might participate in polymerization via its amide (lactam) group; polyamides synthesized in this manner are recognized for their durability, flexibility, and resistance to abrasion and chemicals, finding applications across diverse fields ranging from textiles (such as nylon) to automotive components~\cite{Varghese2022BeyondApplications,Varghese2022BeyondApplications}.

Alkyl substituents can enhance the hydrophobic nature of the molecule, leading us to expect that polymers derived from monomers containing such alkyl groups will exhibit increased hydrophobicity. This change can influence their solubility across various solvents as well as their interaction with water. Such characteristics might be beneficial for applications that demand water resistance or particular solubility properties. Additionally, alkyl groups may impact the mechanical attributes of the polymer. For example, they can serve as plasticizers within the polymer matrix, thereby enhancing the polymer's flexibility. This effect could result in materials that are more flexible or possess improved impact resistance, contingent on the alkyl groups’ chain length and branching.

\subsubsection{Additional benchmarking} We present further benchmarks of the X-LoRA-Gemma model, illustrated in Figure~\ref{fig:perf_MMLU}, utilizing a subset of the MMLU benchmark focused on chemistry, physics, and biology (domains closely aligned with the training data of our model)~\cite{Hendrycks2020MeasuringUnderstanding}. Our findings reveal that both X-LoRA models covered in this paper outperform their respective base models. Specifically, the X-LoRA models surpass both base models (Zephyr model: 54.6\% for the base model compared to 56.6\% for the X-LoRA model; Gemma model: 40.6\% for the base model versus 52.4\% for the X-LoRA-Gemma model). Among them, the X-LoRA-Zephyr model achieves the highest overall performance, although the results of both are fairly comparable.

\section{Conclusion}

Multimodal LLMs have numerous current and prospective applications in scientific fields, alongside a strong demand for highly specialized, domain-specific expertise. Beyond general language processing, the ability to perform computational or generative tasks, as exemplified here with protein mechanics, is essential~\cite{Luu2023BioinspiredLLM:Materials,Buehler2023MechGPTModalities,Buehler2023MeLMProblemsb}. X-LoRA, which is based on open-source language models, allows for the dynamic blending of expertise and knowledge from various domains, as illustrated by several examples, including Figure~\ref{fig:conversation_evolve}. This analysis revealed how the X-LoRA model autonomously and dynamically adapts the use of adapters, combining specific adapters in a sophisticated manner across the LoRA layers. Such intricate mixing of LoRA layers is a consistent pattern observed in all examples examined here, such as in Figures~\ref{fig:QA_weights_sample_0} and \ref{fig:QA_weights}. This observation presents an intriguing challenge to better understand these mechanisms through a physics lens, where an LLM can be viewed as a large-scale graph-forming model leveraging complex graph structures. Consequently, the blending of adaptation experts at deeper layers, as conducted here, modifies these mechanisms to devise effective solutions.

The proposed approach enhances the traditional inference method, which relies on a single forward pass, by incorporating two forward passes. This enables the model to initially ‘consider’ the question and how it might reconfigure itself before producing an answer. This represents a straightforward implementation of ‘self-awareness’, allowing the model to modify its own structure to optimally address a task. Consequently, despite having a relatively modest parameter size of 7 billion, the model can perform reasoning across a broad range of scientific disciplines (biological materials, mathematics, physics, chemistry, logic, mechanics, etc.), substantially boosting its potential for generating novel solutions. Moreover, this approach may inspire the creation of alternative model architectures, such as techniques that generalize the reconfiguration strategy to parts of a model or to entire non-adapted systems. In these scenarios, the scaling head would dynamically reconfigure sections of a trained model or combine multiple models. As demonstrated with X-LoRA, such strategies offer powerful means to extend or unify the capabilities of different models.

We conducted multiple performance evaluations, including a comparison with significantly larger models (Figure~\ref{fig:perf_comp_bioinspiredexam}), where X-LoRA demonstrated clear superiority. The comprehensive comparison presented in Table~\ref{tab:table_qa} revealed that X-LoRA significantly outperforms the base model across a wide range of tasks and application areas. Additional comparisons involved quantitative evaluations of numerical prediction problems, such as protein mechanics and quantum-mechanical molecular properties, where X-LoRA achieved high accuracies with $R_2$ values ranging from 0.85 to 0.96. For context, these results surpass previously published outcomes, for instance, predicting protein unfolding force using the MeLM model~\cite{Buehler2023MeLMProblemsb}, which reported an $R_2$ of 0.78. The ability to accurately handle diverse prediction tasks, inverse design challenges, and complex reasoning was highlighted in the molecular design example (Section~\ref{gemma-section}). This case also demonstrated that X-LoRA can be effectively implemented on various base model architectures (Gemma versus Zephyr/Mistral). Future research could further extend this approach and develop X-LoRA models for larger base models.

A limitation of our method is the requirement for two forward passes, which increases computational cost: one to compute hidden states used as input for the X-LoRA scaling head, and another with $\Lambda$ applied to the adapters. Nevertheless, the first pass benefits the X-LoRA scaling head by leveraging the LLM's intrinsic capabilities, unlike learning the knowledge independently by training a larger scaling head. For efficient deployment in model-serving environments, it is necessary to maintain separate key-value caches for both the scaling and forward passes. To enhance inference speed, we created \texttt{Mistral.rs}, a LLM serving platform implemented in Rust~\cite{matsakis2014rust} that supports X-LoRA and incorporates various performance optimizations. Despite these challenges, a major advantage of our approach is its compatibility with any model without requiring modifications to the internal architecture. We consider this a significant benefit, particularly as it enables leveraging the extensive resources available within the Hugging Face ecosystem.

One benefit of the proposed method is that it offers an easy implementation path for any existing LLM (for example, with the release of the X-LoRA code, it can be readily applied to any model within the Hugging Face ecosystem). Another benefit is that the scaling head leverages the inherent capabilities of the utilized LLM and learns advanced techniques for optimally combining the adapters to generate particular responses. During the training phase of X-LoRA, it is ideal to supply complex examples of question-answer pairs or conversations that enable the model to discover the most effective combinatorial strategy. We conducted several analyses related to protein design and mechanics. For example, Figure~\ref{fig:MJB_20} and Figure~\ref{fig:MJB_21} showcase a stability assessment of multiple protein sequences, including comparative evaluations, reasoning over the predicted outcomes, and a mechanistic explanation that provides a fundamental chemical basis for the predicted behavior. As illustrated in Fig.~\ref{fig:MJB_21}, this was directly validated through a structural analysis of the protein’s folded 3D conformation.

Regarding potential future work, further studies could investigate a wider variety of adapter experts, expanding upon the initial expertise included in our LoRA set. The X-LoRA model outperforms the base model alone or the base model combined with a single adapter, offering answers to intricate questions by leveraging specialized knowledge from the collection of adapters. There also appear to be interactions between the various adapters, as indicated by the complex mixing effects of scaling weights across different layers. This phenomenon could be explored more deeply and compared with approaches such as SLERP, which combines models using alternative strategies. While we trained the X-LoRA scaling head using a subset of training data, consisting of a few hundred samples from each of the original datasets used for adapter development, a more targeted training set designed specifically for this purpose could be created. For instance, we observed that multiple-choice questions tend to activate adapters focused on reasoning because these were trained on such datasets. Additional prompting strategies are needed to assist the scaling head in grasping particular domains relevant to the questions, for example, through specialized system messages or other techniques. Overall, the creation of suitable training sets is an area warranting further investigation, including methods to effectively induce mixing of layer-wise scaling. Our experiments tracking these dynamics over extended conversations demonstrate that the model can clearly invoke varying scaling mechanisms to best address specific tasks, which is an encouraging result.

We identified intriguing patterns in how experts are activated across the layers, with several noteworthy observations elaborated upon in the paper. This analysis has the potential for further expansion and could provide valuable insights into the mechanisms and benefits of mixing different model components or sets of layers. Such an exploration would be especially compelling when the X-LoRA technique is applied beyond protein mechanics and protein design, as demonstrated here. We reserve this for future research.

Another significant direction involves deeper exploration of agentic modeling, as exemplified in this study through the use of two adversarially interacting X-LoRA models. Future research could incorporate more than two models to foster richer interactions and introduce additional functionalities that push models into novel generative domains. This approach may also facilitate the development of methodologies that incorporate physics-based modeling or other validation protocols, utilizing code-writing/executing agents or agents leveraging function calling or other processing methods~\cite{Ghafarollahi2024ProtAgents:Learning}.

\section{Materials and Methods}

The X-LoRA code and illustrative examples can be accessed at \url{https://github.com/EricLBuehler/xlora}. Model weights and datasets related to the X-LoRA described in this manuscript, together with further examples, are hosted at \url{https://huggingface.co/lamm-mit/x-lora}. The X-LoRA-Gemma model is available at \url{https://huggingface.co/lamm-mit/x-lora-gemma-7b}. The \texttt{Mistral.rs} inference engine is downloadable from \url{https://github.com/EricLBuehler/mistral.rs}.

\subsection{Mathematical formulation of X-LoRA} We define the $\mathbf{scalings}$ as a matrix comprising token- and layer-level coefficients for LoRA adapters. The X-LoRA $\mathbf{scaling}$ $\mathbf{head}$ receives hidden states from the base model and applies linear fully-connected layers to predict these scalings. A softmax function is subsequently applied along the final dimension to ensure the scalings across all experts sum to one. ReLU activation functions and dropout layers are interleaved between linear fully-connected layers within the X-LoRA scaling head. The size of the hidden layers is configurable, allowing for effective dimensionality reduction from the hidden dimension to the scaling dimension or enabling a widening-narrowing scheme analogous to the feed-forward layer in a transformer block.

The scalings form a matrix $\Lambda \in \mathbb{R}^{s \times l \times n}$, where $s$, $l$, and $n$ represent sequence length, number of layers, and number of adapters, respectively. For each layer, the relevant scalings are extracted as a matrix $\lambda \in \mathbb{R}^{s \times n}$. These scalings $\lambda$ are then applied to each LoRA adapter by broadcasting the corresponding X-LoRA scaling $\lambda_i \in \mathbb{R}^{1 \times 1 \times s}$ to the inputs of the decomposition matrices $A$ and $B$.

Given $W_0 \in \mathbb{R}^{d \times k}$, $B_i \in \mathbb{R}^{d \times r_i}$, and $A_i \in \mathbb{R}^{r_i \times d}$ with rank $r_i \ll {\rm min}(d,k)$ unique to each adapter $i$, the forward pass of an X-LoRA layer is described by the equation:

\begin{equation}
    h = W_0 x + \sum_{i=1}^{n} B_i A_i (x \cdot \lambda_i) \alpha_i
\end{equation}

Representations of the scalings are shown in the main text. These visualizations were produced either via bar plots or with Matplotlib’s \texttt{contourf} function~\cite{MatplotlibDocumentation}.

\subsection{X-LoRA implementation algorithm and code} The implementation selected here prioritizes user-friendliness and the flexibility to seamlessly adapt the method to a wide range of models, particularly those within the Hugging Face ecosystem. The X-LoRA algorithm functions by:

\begin{enumerate}
    \item Keeping the base model and adapters fixed, then optimizing performance.
    \item Combining LoRA adapters based on the scaling predictions produced by the X-LoRA scaling head.
\end{enumerate}

The entire codebase is developed using Python and PyTorch~\cite{Paszke2019PyTorchLibrary}.

\subsubsection{Dual forward pass strategy for self-aware inference} Since X-LoRA necessitates computing hidden states to predict scaling factors, we adopt a dual forward pass approach: the initial pass allows the model to assess the task and determine how to adjust itself, followed by the actual inference pass. This mechanism introduces a form of `self-awareness', enabling X-LoRA to allocate extra computation time for introspection instead of immediate output generation.

We refer to the first pass as the scaling pass and the second as the forward pass. During the scaling pass, the scalings $\Lambda$ are initialized to zero. We also experimented with initializing $\Lambda$ to ${n}^{-1}$ and observed comparable outcomes, signifying robust performance (this parameter can be configured within the X-LoRA package). To facilitate straightforward application of X-LoRA across diverse models, we implement a forward pass hook. This approach is crucial for linking the scaling and forward passes and represents the central feature enabling compatibility with all models.

\begin{enumerate}
    \item $\mathbf{Scaling}$ $\mathbf{pass}$: Conducted by running the base model with adapters fixed at constant values. The resulting hidden states feed into the X-LoRA scaling head to compute $\Lambda$.
    \item $\mathbf{Forward}$ $\mathbf{pass}$: Utilize the scaling factors $\Lambda$ to combine the adapters during the model’s forward pass. The output logits are then returned as the X-LoRA model’s prediction.
\end{enumerate}

It is important to note that sharing the same key-value cache between the scaling and forward passes can disrupt accurate inputs to the scaling head because the key-value cache persists and accumulates data across both passes. Therefore, we disable the key-value cache during training and inference and have implemented appropriate management of this in \texttt{Mistral.rs}.

\subsubsection{Freezing weights and enabling model segments for training} The base model and adapters are kept frozen ({\it i.e.}, their weights remain fixed) so that the X-LoRA scaling head is the sole component trained in an X-LoRA model. However, the presented code allows the option to unfreeze all adapters in addition to the X-LoRA scaling head. Moreover, one may choose to train the entire model simultaneously—comprising the X-LoRA scaling head, adapters, and the underlying foundation model—potentially using distinct learning rates. Alternatively, future developments based on the concepts introduced here could facilitate creating scaling functions that operate between multiple foundation models, enabling a transparent implementation of SLERP-type modeling. Such directions could be pursued in subsequent research.

\subsubsection{X-LoRA layers} An X-LoRA layer functions by scaling and ultimately aggregating outputs from multiple adapters. Each X-LoRA layer encapsulates the original LoRA layer, thus retaining access to its corresponding adapters.

To effectively apply X-LoRA across various models, the X-LoRA approach replaces the original LoRA layers with newly constructed X-LoRA layers, then integrates their forward methods into the LoRA layers. This process leverages Python’s object-oriented capabilities to facilitate this swapping mechanism.

In brief, during the forward pass, each X-LoRA layer performs the following operations (refer to Figure \ref{fig:general_arch}):

\begin{enumerate}
    \item $\mathbf{Scaling}$: Scale the input signal of each LoRA adapter by its respective scaling factor $\lambda_i$.
    \item $\mathbf{Forward}$: Execute the forward pass through the encapsulated LoRA layer.
\end{enumerate}

\subsection{Efficient Inference using \texttt{Mistral.rs} implemented in Rust} To enhance inference efficiency, we developed \texttt{Mistral.rs}, a Rust-based LLM serving platform that implements X-LoRA~\cite{matsakis2014rust} alongside several optimizations to boost performance. In particular, the following techniques are applied, listed here without any implied order:

\begin{itemize}
    \item LoRA adapter weight stacking: When all $A_i$ and $B_i$ matrices share identical dimensions across LoRA adapters, it becomes possible to stack these $A_i$ and $B_i$ weight matrices along their first dimension.

Mathematically, given $W_0 \in \mathbb{R}^{d \times k}$, for adapter matrices $A_i \in \mathbb{R}^{r \times d}$ and $B_i \in \mathbb{R}^{d \times r}$ with $r_i \ll {\rm min} (d,k)$, we concatenate these matrices such that $A\in \mathbb{R}^{i \times r \times d}$ and $B\in \mathbb{R}^{i \times d \times r}$. The parameters $\alpha_i$ are also applied during the initialization of the stacked matrices.

To implement scalings, we rearrange the dimensions so that $\lambda \in \mathbb{R} ^ {i \times 1 \times 1}$. These scalings are then applied via broadcast multiplication.

\item Non-granular scalings To decrease the frequency of running the scaling pass, we introduce an option for non-granular scalings. This approach caches the scalings after $n$ completion runs, taking advantage of the KV cache which maintains the sequence length at one for the remainder of the request. Utilizing this method significantly lowers latency by reducing the number of times the base model needs to be invoked.

\item Fused Compute Unified Device Architecture (CUDA) kernels For optimizing the forward pass, we employ fused CUDA kernels. These fused kernels allow combining complex operations, such as Rotary Embedding (\cite{su2021roformer}), which would normally be executed in sequence, into a single kernel that can run in parallel on a GPU. Our kernels are based on the vLLM implementation (\cite{kwon2023efficient}).

\item Quantization The inference engine supports quantization, enabling the use of a quantized base model. Initial experiments demonstrate that quantization down to 8 bits performs well even for models originally trained with \texttt{bfloat16} precision, for example.

\end{itemize}

\subsection{Datasets}

Each adapter is trained using specialized datasets, summarized in Table~\ref{tab:table_loradata} along with citations to the sources and original literature.

\subsection{Training strategy} Training proceeds in stages. We start with the Zephyr-7B-$\beta$ model~\cite{HuggingFaceH4/zephyr-7b-betaFace}, which was developed on top of the Mistral-7B model \cite{Jiang2023Mistral7B}, and train a series of adapters using the datasets described above. The Zephyr-7B-$\beta$ tokenizer is used as well.

The first eight adapters are trained using question-answer pairs provided by the datasets (refer to Table~\ref{tab:table_loradata}).

The protein mechanics tasks are trained similarly, but employ specific forward and inverse instruction sets. The bidirectional instruction tasks include:

\begin{quote}
\tiny {\texttt{CalculateForce< G E E C D C G S P ....> [0.310]\\
CalculateEnergy< G E E C D C G S P ....> [0.220]\\ 
CalculateForceHistory< G E E C D C G S P ....> [0.004,0.034,0.125,0.142,0.159,0.102,0.079,0.073,0.131, ...]\\
GenerateForce<0.310>\,[ G E E C D C G S P ....]\\
GenerateEnergy<0.220>\,[ G E E C D C G S P ....]\\
GenerateForceHistory<0.004,0.034,0.125,0.142,0.159,0.102,0.079,0.073,0.131, ...>\,[ G E E C D C G S P ....] }}
\end{quote}

These correspond to three distinct tasks: computing the maximum unfolding force of a protein, determining the unfolding energy, and capturing the force-deformation history. Each of these three forward tasks is paired with a corresponding inverse task of the same type, which produces a protein sequence as output, resulting in six tasks in total.

The rank of each LoRA adapter is set to $r=16$ with a scaling factor $\alpha=16$. Training for each adapter involves five specific target modules: \texttt{v\_proj}, \texttt{k\_proj}, \texttt{q\_proj}, \texttt{gate\_proj}, and \texttt{o\_proj}, applying a dropout rate of 0.05. Typically, LoRA adapters are trained using low ranks ranging from 4 to 32. This is because prior studies have indicated that increasing the LoRA adapter rank seldom enhances performance, while it does raise computational costs, as lower ranks yield more compact models with fewer parameters~\cite{hu2021lora}. The parameters $r$, $\alpha$, and others are chosen based on widely adopted values in existing literature known for their efficacy. Additionally, previous research conducted by the author has tested various configurations and concluded that the values used here represent a sensible selection~\cite{Luu2023BioinspiredLLM:Materials}.

The X-LoRA scaling head is trained using approximately 2,000 samples, which represent a subset of the original dataset employed to train the adapters (additionally, we augmented the training data with GPT-4 generated multiple choice question samples). We utilize a \texttt{paged\_adamw\_8bit} optimizer, applying gradient clipping at 0.3, a learning rate of $2\times 10^{-4}$ with warmup, and perform four gradient accumulation steps, with a batch size set to one. The X-LoRA model is trained for roughly 10,000 steps, reaching a loss near 0.28 at the end of training.

Given that the basic Mistral architecture underlying the foundation model consists of 32 layers, and adapters are applied to 5 transformer layers each, this amounts to a total of 160 LoRA layers, all of which are replaced by X-LoRA layers. The X-LoRA scaling head predicts scaling factors for each of these 160 LoRA layers; for nine LoRA experts, this corresponds to $160 \times 9 = 1440$ $\lambda_i$ values computed per token. We implement a single linear layer with ReLU activation and dropout at $p=0.2$, comprising around 5 million parameters. The codebase is designed with flexibility in mind, allowing the construction and investigation of more complex deep classifier heads for other use cases.

Both LoRA adapters and the X-LoRA scaling head are trained through next-token prediction using cross-entropy loss (we determine the likelihood assigned by the model to the correct subsequent token in a sequence given the preceding tokens; the objective during training is to maximize the probability of the accurate next token as indicated in the training data). We implement token shifting, whereby the sequence is offset by one token to generate the target sequence for the model to predict. For instance, if the input sequence is ``Proteins are fundamental to", the corresponding training target sequence would be ``are fundamental to biology". Hence, the model attempts to predict the following token in the sequence.

As depicted in Figure~\ref{fig:hierarchical_concept}, during the training of the X-LoRA scaling head, only the scaling head’s parameters are updated while all other parameters remain fixed. The original Zephyr tokenizer and prompt template are employed.

\subsection{X-LoRA-Gemma model} The X-LoRA-Gemma model is constructed similarly to the X-LoRA model described above, but it is built upon the Gemma-7B-it model~\cite{Google/gemma-7b-itFace}, utilizing four adapters (refer to the main text for more details). These four adapters are trained on datasets covering mechanics and materials, protein mechanics, bioinspired materials, and an additional quantum-mechanics molecular properties dataset QM9 (see Table~\ref{tab:table_QM9})~\cite{Ruddigkeit2012EnumerationGDB-17,Ramakrishnan2015ElectronicSpace}.

Each LoRA adapter has rank $r=16$ with $\alpha=16$. The adapters are trained targeting seven modules: \texttt{q\_proj}, \texttt{o\_proj}, \texttt{k\_proj}, \texttt{v\_proj}, \texttt{gate\_proj}, \texttt{up\_proj}, and \texttt{down\_proj}, applying a dropout rate of 0.05. The rank for each LoRA adapter is selected as $r=16$ with $\alpha=16$. Considering four adapters and seven target modules, and given there are 28 layers in the Gemma-7B base model, a total of $28 \times 28 = 784$ $\lambda_i$ values are computed per token.

The initial two adapters (bioinspired and mechanics of materials) incorporated into this model are trained on question-answer pairs as supplied by the datasets (see Table~\ref{tab:table_loradata}).

Protein mechanics tasks are trained similarly to the approach outlined earlier, employing bidirectional forward and inverse instruction sets. Furthermore, an adapter is trained using the QM9 dataset, involving these two tasks:

\begin{quote}
\tiny {\texttt{CalculateMolecularProperties<CC(C)N1CC1C=O> [0.098,0.358,0.581,0.395,0.309,0.330,0.570,0.649,0.519,0.519,0.519,0.518]\\
GenerateMolecularProperties<0.098,0.358,0.581,0.395,0.309,0.330,0.570,0.649,0.519,0.519,0.519,0.518> [CC(C)N1CC1C=O]\\ 
...
}}
\end{quote}

These correspond to a total of two tasks: predicting molecular properties (the forward task) and the inverse task of generating a molecule with specified target properties. All 12 properties included in the QM9 dataset are either predicted or designed for, respectively (refer to Table~\ref{tab:table_QM9}).  
This model contains four adapters in total.  

The X-LoRA scaling head for this model comprises three layers, with an intermediate layer size of $3072 \times {FF}_{\rm fac} = 12,288$, where ${FF}_{\rm fac} = 4$ serves as a feed-forward multiplier. We implement linear layers with ReLU activation and a dropout rate of $p=0.2$, resulting in approximately 47 million parameters.  

The X-LoRA-Gemma scaling head is trained on roughly 17,000 samples. These samples were selected as a subset from the training dataset used for adapter training; it is important to note that although the X-LoRA-Gemma model includes only four adapters, the samples were drawn from the entire training dataset utilized for the model described previously, supplemented with samples from the QM9 dataset. This approach was taken to explore whether the X-LoRA model could acquire additional capabilities during the scaling head training phase. Given that the loss converged well to 0.58 by the end of training, alongside strong performance even on tasks outside the explicit training scope (e.g., chemistry, as illustrated in the molecular design example presented in the paper), we conclude that this methodology generally produces favorable outcomes.

We employed a \texttt{paged\_adamw\_8bit} optimizer with gradient clipping set to 0.3, a learning rate of $1\times 10^{-4}$ including warmup, and accumulated gradients over four steps, using a batch size of one. The X-LoRA-Gemma model was trained for approximately 62,000 iterations.

We utilized the native Gemma tokenizer along with its prompt template.

\subsection{Adversarial agentic modeling} Our adversarial agentic framework involves creating two X-LoRA agents. One agent specializes in posing questions, initiating with an initial query and continuously generating further questions aimed at uncovering flaws and challenges in the answers. The other agent serves as the responder, providing ongoing replies. This setup is implemented using \{Guidance\}~\cite{Guidance-ai/guidance:Models.}.

For the example linking music and mechanics mentioned in the text, the characteristics of the two agents are defined as follows. First, the physicist and questioner:

\begin{quote}
\small\texttt{You are a critical physicist. You participate in a discussion from the physics perspective. Keep your responses concise and consistently challenge statements provocatively.

As a creative thinker, you introduce ideas from other disciplines and push conceptual boundaries.}
\end{quote}

Second, a philosopher (in certain examples this role is adapted to represent a biology expert):

\begin{quote}
\small\texttt{You are a philosopher with expertise in biology, chemistry, and mathematics. You engage in the discussion.

Provide brief yet precise and imaginative answers. You generate excellent ideas and novel lines of reasoning, always maintaining logical coherence.}
\end{quote}

The question-asking agent is explicitly instructed to reflect on the question and preceding responses, crafting new insightful queries. This is facilitated through a prompt formulated as follows:

\begin{quote}
\small\texttt{Examine this question and answer.

\#\#\# Question: <question>

\#\#\# Response: <response>

\#\#\# Instruction: Provide a SINGLE follow-up question that critically examines the response.

Do NOT provide an answer or commentary on the question at this stage.

The single question is: }
\end{quote}

After the dialogue has continued for $N$ turns, the model is assigned to 1) summarize the primary insights discussed, 2) enumerate the key insights as bullet points, and 3) determine the most significant conclusion.

The raw conversation text serves as the foundation for subsequent analysis through graph creation.

\subsection{Knowledge graph generation} We employ Zephyr-7B-$\beta$ to extract triplets from the text, following the method outlined in~\cite{Rahulnyk/knowledge_graph:QnA} with additional enhancements inspired by the Llama Index graph generation algorithm~\cite{Liu_LlamaIndex_2022}. Graphs are rendered using NetworX\cite{Networkx/networkx:Python} and Pyvis~\cite{WestHealth/pyvis:Graphs.}. A key instruction provided is, consistent with the method proposed in~\cite{Liu_LlamaIndex_2022}:

\begin{quote}
\small\texttt{Your task is to derive the ontology of terms appearing in the provided context. These terms should capture the main concepts according to the context.

Present your output as a JSON list. Each item in the list includes a pair of terms and their relationship, structured as follows: [...] }
\end{quote}

We guide the model to construct triplets of nodes and edge attributes as follows, adopting the technique used in~\cite{Liu_LlamaIndex_2022}:

\begin{quote}
\small\texttt{\begin{itemize}
            \item "node\_1": "A concept extracted from the ontology"
            \item "node\_2": "A related concept extracted from the ontology"
            \item "edge": "A concise description of the relationship connecting node\_1 and node\_2"
        \end{itemize}
        }
\end{quote}

We also supply the model with examples of ontological analysis, for instance:

\begin{quote}
\small\texttt{
       Context: ```Spider silk is a strong natural fiber used to catch prey in a web.``` }
\end{quote}

We then provide example triplets, such as:

\begin{quote}
\small\texttt{\begin{itemize}
            \item "node\_1": "spider silk"
            \item "node\_2": "fiber"
            \item "edge": "is"
        \end{itemize}
        }
\end{quote}

The assembled graph data is then organized into communities using the Girvan-Newman algorithm~\cite{Girvan2002CommunityNetworks}.

\subsection{Visualization of molecular structures}

We utilize PyMOL~\cite{Schrodinger2015The1.8} to visualize and examine the predicted protein conformations. Various scripts and functions within the software are applied to identify specific protein characteristics, such as secondary structure, hydrophobic/hydrophilic regions, disulfide linkages, and hydrogen bonding.

\section*{Data Availability Statement} The codes and datasets underpinning the results of this research are freely accessible at \url{https://github.com/EricLBuehler/xlora} and \url{https://huggingface.co/lamm-mit/x-lora}. Additional data sources used in this work are cited (refer to Table~\ref{tab:table_loradata} for details).

\end{document}